\section{Entry from a positive Gaussian coefficient}
\label{en:section}

Fix an admissible finite interaction set $J$.  We use
$A_0=-\Delta_J/v_J^2$, $k=\beta/v_J^2$, and the structural coefficient
$a(k)$.  In this section $k_h$ is fixed and $a_h=a(k_h)>0$.
In particular $a_h\geq8\pi$.  The input temperature $k$ and the
reference coefficient $a'$ will vary near $k_h$ and $a_h$, respectively;
$\alpha=\sqrt{a'}$.  A derivative in $k$ always holds $a'$ fixed.
Fix a compact interval $I\Subset(0,\infty)$ with
$a_h\in\operatorname{int}I$, and take $a'\in I$ in all comparison-energy
estimates below.
The results on mixed domains that are used below are
Propositions~\ref{an:mixed-scaling}, \ref{an:pin-interaction},
\ref{an:conditional-pins}, and \ref{an:gaussian-localization}.
The microscopic scaling hypotheses are imposed only at $k_h$.
The pure real-Gaussian locality bounds in
Proposition~\ref{an:gaussian-localization} are uniform on $I$;
each comparison of topologies uses the same $a'$.

\subsection{Statement and the observation identity}

For a dyadic integer $s$, let $B_s$ average over the disjoint square
cells of side $s$, and introduce the observations
\begin{equation}\label{en:observation}
                 Y=B_sh+\eta,\qquad \eta\sim N(0,I).
\end{equation}
The noises are independent of $h$.  Work modulo common translations by
$2\pi$, with the common mean in one period.  The comparison observation
is $B_s\alpha\psi+\eta$, where the mean-free covariance of $\psi$ is
$A_0^+$ and its common mean is uniform modulo $2\pi/\alpha$.
Both conventions give ordinary Gaussian formulas on tests annihilating
constants.  On a finite torus the observation density of the height
model is proportional to $Z_{\rm full}(y;k)$, in the notation of
Section~\ref{en:patches} below and of Proposition~\ref{an:mixed-scaling}.
Its Gaussian comparison density is proportional to
$\exp(-E_{\rm full}(y)/2)$.  Thus its exact Radon--Nikodym derivative is
\begin{equation}\label{en:exact-density}
 \frac{Z_{\rm full}(y;k)e^{E_{\rm full}(y)/2}}
 {\displaystyle\int Z_{\rm full}(z;k)e^{E_{\rm full}(z)/2}
                      \,d\gamma_{a',s}(z)}.
\end{equation}
This identity introduces observations of the original heights; it does
not change their interaction.

Here and below the block supports, point norm, and ordinary regulator
$W_u^\kappa$ have the definitions of Lemmas~\ref{rg:geometry}
and \ref{rg:regulator}.  Enlarging their fixed collars to contain the
finite range of $J$ is understood.  We will multiply this regulator by
a nonnegative quadratic factor $e^{V_j}$, defined in
Equation~\eqref{en:variational}.  Write
\begin{equation}\label{en:norm}
 \|F\|_{u,A,V}=\sup_{X,\psi}
 A^{|X|}\frac{|F(X)|_{u,X,\psi}}
 {W_u^\kappa(X,\psi)e^{V_j(X,\psi)}},\qquad u=mL^j.
\end{equation}
The point norm includes the full analytic derivative sum at radius
$h_0$; it is not a norm restricted to bounded fields.

\begin{theorem}[Local entry at every positive coefficient]\label{en:entry}
Fix $J,k_h,a_h$ as above.  Fix admissible parameters for the small-activity
maps of Proposition~\ref{wh:map}, including their block factor $L$, activity
weights $A$, and analytic radius $h_0$, with the regulator choices allowed
by Lemma~\ref{en:regulator}.  For every $\epsilon>0$ there are a fixed
dyadic integer $q$, an arbitrarily large dyadic $R$, $P=R^5$, a dyadic
$s$, and a neighborhood
$\mathcal U\subset(0,\infty)\times\operatorname{int}I$ of $(k_h,a_h)$,
such that the following assertions hold at $m=qRs$.

On the plane there are real, even, square-covariant local activities
$K_m(X;k,a')$, invariant under translations of the block grids and
under common field translation by $2\pi/\alpha$, with
\begin{equation}\label{en:small}
                 \|K_m(k,a')\|_{m,A,V}\leq\epsilon
                 \quad ((k,a')\in\mathcal U).
\end{equation}
Their subset expansion on every dyadic torus fitting the hierarchy and
with side $n\geq sP$ is an exact representation of
Equation~\eqref{en:exact-density}, after integrating the observation
noise and the Gaussian covariance below side $m$, up to a
field-independent factor that cancels upon normalization.  The torus
activities obey the same estimate.  A support whose padded neighborhood
unwraps has exactly its plane activity.  This copying property persists
under the local maps up to their specified stopping scale.
Here fitting means that $n/s$ is a multiple of $qR$, the torus uses
the same grid origin, and $n/(mL^j)$ is an integer for every retained
coarse grid.

On a smaller neighborhood the activities are analytic in $k$ at fixed
$a'$.  For every integer $r\geq0$ there is $C_r<\infty$ such that
\begin{equation}\label{en:derivatives}
 \|\partial_k^rK_m(k,a')\|_{m,A,V}\leq C_r\epsilon.
\end{equation}
The constants may depend on $s,R,J,k_h$ and the chosen neighborhoods.
Finite-support evaluations, their field derivatives, and scalar
localization projections are continuous in $a'$.

The choices are made in the following order: the downstream norm
parameters and the imaginary field strip widths are fixed first;
then $R$ and $P$; then a finite large-field threshold; then $s$; and
finally $\mathcal U$.  No uniform scaling assertion at temperatures
other than $k_h$ is part of either the hypotheses or the proof.
\end{theorem}

One may start the split map with $t=z=0$ and $\mathcal R=K_m$.
Its first localization, or its balanced coordinates, consequently has
size $O(\epsilon)$.  The entry estimates hold on the full positive
neighborhood $\mathcal U$.  The near-critical split continuation is
used on its subset with
$k_0\le a'\le k_0+\delta_{\rm ref}$, the one-sided interval of
Proposition~\ref{wh:map}.  When $a_h=8\pi$, this includes the nearby
warmer true references used in the endpoint argument.  On a compact
reference interval strictly above $8\pi$ one may instead use the
gapped continuation.

The remaining subsections prove Theorem~\ref{en:entry}.
The isolation expansion in Section~\ref{en:patches} produces local
differences controlled on bounded real observations by the mixed-pin
results; the observation-energy regulator of Lemma~\ref{en:regulator}
and the real pattern estimates then control large observations.
Gaussian averaging, complex comparison, and analytic interpolation
upgrade these bounds to the full norm in Equation~\eqref{en:norm},
and connected pattern sums complete the construction.

\subsection{A finite, covariant isolation expansion}\label{en:patches}

Use the tiles specified in Proposition~\ref{an:mixed-scaling} and shown in
Figure~\ref{an:mixed-geometry-figure}: in units of
$s$ cells start with the square grid of side $R$, remove the square
of side $R/2$ centered at each vertex from its incident squares, and
declare each removed square a tile as well.  The retained parts of the
original squares and these vertex squares partition the cells.  Cuts
are placed between site layers.  Thus every junction has at most three
incident tiles and every pair meeting there also shares a segment.
The mixed-domain results apply to unions of these tiles, to graphs with
selected tiles isolated by deletion of all their exterior $J$-bonds,
and to their $R/4$-aligned box clippings.  Components are defined before
any hard pins, and point contacts do not identify components.

Here is a completely finite ordering of the bond deletions.  Color tile
centers by their tile type and the square-symmetry orbit of their
coordinates modulo a dyadic period $q$, in units of $R$.  Within each
color, group centers into proximity components, with proximity radius
$100$ initially.  Choose $q$ greater than $16$ times this radius and
increase both fixed constants if collars require it.  Each group
contains at most eight centers: a shortest path repeating a residue
label has at most eight steps, hence displacement less than $q$;
it must therefore repeat the same center.  Removing such a repetition
leaves at most the eight distinct orbit labels.  Distinct groups of
one color have disjoint enlarged neighborhoods.  Order the finitely
many colors once and for all.  An event isolates every member of its
group, also from the other members, counting a bond only if it is
still present.

For each group its core contains the tiles at center distance at most
$1$ from a member; its open patch contains those at distance at most
$5$; its guard contains the patch tiles at distance at least $3$.
All distances here are sup distances in units of $R$.  For sufficiently
large $s$, every changed bond is inside the core and the guard separates
the core from the exterior.  When a guard itself is tested as a pin,
use the enlarged patch through distance $12$, with outer guard starting
at distance $10$.  These constants may be enlarged a fixed number of
times.  In all subsequent uses the proximity radius is enlarged first
so that simultaneous marked interiors are separated.  Grid origins
are chosen at elementary-square centers; square symmetries and block
translations preserve all these rules.

For an event $i$, let $o$ and $c$ be its current open patch before and
after deletion of its designated bonds.  Put
\begin{equation}\label{en:replacement}
 B_i(y)=C_i\exp\{-\tfrac12(E_o(y)-E_c(y))\},
 \qquad C_i=\frac{Z_o(0;k_h)}{Z_c(0;k_h)}.
\end{equation}
The energies use the current $a'$, whereas the scalar constants always
use $k_h$.  In the height sum replace the old bond weight by $B_i$ times
the cut weight plus their difference.  Selecting the difference marks
the event.  After a mark skip every subsequent event with a member
within distance $30$ of that marked group, increasing this constant
with the collars if necessary.  All events in a group are skipped
together.  Events in one color are processed simultaneously.

This defines an exact binary expansion: at each considered event the
two summands add to the original integrand.  A later considered patch
does not meet an earlier marked core or guard, so its replacement is
independent of which of the two signed terms in that earlier difference
is being evaluated.  Marked interiors have no mutual bonds.  Once all
batches have been processed, any undeleted cross-tile bond has an
endpoint in a fixed halo of a mark.  Conversely every difference from
the path with no marks is confined to such halos.  These facts follow
directly from the skip distance and the bounded number of batches;
the necessary halo radius is a fixed constant in $R$ units.

Let $B_i^{(0)}$ denote the factors along the path with no marks.  Divide
the full expansion by
\begin{equation}\label{en:base}
       \prod_iB_i^{(0)}(y)
       \prod_b Z_b(0;k_h)e^{-E_b(y)/2}.
\end{equation}
Here $b$ runs through the isolated tiles.  A halo of a mark retains its
whole signed height sum and every affected replacement factor.
Outside all halos, expand each remaining factor as $1+F_b$, where
\begin{equation}\label{en:tile-error}
            F_b(y)=\frac{Z_b(y;k)e^{E_b(y)/2}}{Z_b(0;k_h)}-1.
\end{equation}
Selecting $F_b$ is called selecting an additional tile.  All factors
whose rules and inputs are identical to the path with no marks cancel
exactly; in particular no comparison of centered normalizers over the
unaffected volume is made.

Let $E_0$ be the quadratic energy in Equation~\eqref{en:base}.
Work first on the plane.  To compare a whole batch, telescope its
events one at a time in a finite exhaustion.  The intermediate
topologies are permitted by the real-Gaussian locality statement
of Proposition~\ref{an:gaussian-localization}, uniformly for $a'\in I$.
The enlarged patches of distinct groups in that color are disjoint,
so the open-patch proxy of an event is the one in the simultaneous
rule.  Its true-minus-proxy error retains exponential decay both
from its core and in the separation of the two observation indices,
as well as the guard factor $e^{-cR}$.  Summing over the event
locations is finite with this retained decay.  The same estimates
pass from the finite exhaustion to the whole batch.  Summing the
finitely many batches gives a matrix $M$ for $E_{\rm full}-E_0$ with
zero row sums and
\begin{equation}\label{en:link-matrix}
 |M_{xy}|\leq C e^{-cR-c|x-y|}.
\end{equation}
Cell indices are used here.  The constants are uniform on $I$.

We verify the exact image formula on the allowed tori.  Every
no-mark batch endpoint is invariant under translations by $qR$ cells.
Put $N=n/s$.  When $N$ is a multiple of $qR$ and the grid origins
agree, this pattern descends to the torus.  Every group and proxy
patch has diameter at most a fixed multiple of $R$; after $R$ is
large, $N\ge P=R^5$ makes projection injective on each such patch.
For any no-mark batch endpoint $\tau$, its plane observation
precision is
\[
 T_{\tau,a'}=A_{0,\tau}/a'+B_s^*B_s .
\]
Cell Poincar\'e and the unit squared-average mass give coercivity,
uniformly for $a'\in I$, and the source-tail bound in
Proposition~\ref{an:gaussian-localization} gives exponential cell
decay to the inverse column with source $B_s^*e_y$, where $e_y$ is
the unit cell vector.
The sum of its translates with $y$ replaced by
$y+z$, $z\in N\Z^2$, is locally absolutely convergent.  Applying
the finite-range precision termwise gives the periodized source.
Uniqueness from the same coercivity on the torus identifies this
sum with its torus inverse column.  Thus for the comparison-energy
matrices $\mathsf E_\tau=\Id-B_sT_{\tau,a'}^{-1}B_s^*$,
\[
 \mathsf E_\tau^{(n)}(x,y)
   =\sum_{z\in N\Z^2}\mathsf E_\tau^{(\infty)}(x,y+z).
\]
The proxy matrices, extended by zero outside their patches, have
finite support.  Their torus copies are likewise the sums over
the corresponding event lifts, because those patches unwrap with
the same phase.  Subtracting the before-and-after energy matrices
and these proxies proves
$M^{(n)}_{xy}=\sum_{z\in N\Z^2}M^{(\infty)}_{x,y+z}$.
This uses only the periodic batch endpoints; the intermediate
topologies in the preceding finite telescope need not be periodic.
Write the zero-row-sum
quadratic form as a sum of weighted $(y_x-y_{y+z})^2$, retaining each
image displacement $z$ separately, and expand
\begin{equation}\label{en:links}
       e^{(E_{\rm full}-E_0)/2}
       =\prod_{(x,y,z)}\{1+\ell_{xy,z}(y)\}.
\end{equation}
The definition of $\ell$ is the exponential of the corresponding
quadratic summand minus one.  Its absolute estimates have cost
$Ce^{-c(R+d)}$, where $d$ is the full displacement in cell units,
times a small positive quadratic cost along a path connecting its
endpoints.  Indeed $|e^t-1|\leq |t|e^{|t|}$ and
$(y_x-y_{y+z})^2\leq d\sum_{e\text{ on path}}(\nabla_e y)^2$;
polynomial factors in $d$ can be absorbed by weakening $c$.
Include both coordinate-order paths to preserve square symmetry.
The sum of all reserved path costs is exponentially small per edge,
because the number of possible endpoints within path length $d$ is
polynomial in $d$.  Equation~\eqref{en:links} is absolutely convergent
at a fixed finite-torus field.  Its integrated absolute convergence
will also follow from Equation~\eqref{en:averaged-pattern} below.

Cover every mark and selected tile by its required tile halo, padded
also for the estimates below.  Cover a link by its complete path,
including its winding if it projects to a torus.  Take all $m$-blocks
met by these covers.  Apply the polynomial positivity and range argument
for the fixed-$J$ entry family in Lemma~\ref{wh:kernels}.
Since $A_0=Q/v_J^2$, the covariance
\begin{equation}\label{en:entry-covariance}
 C_{<m}=C^{\rm en}_{<m}=A_0^{-1}(1-\mathsf P_m)
\end{equation}
is $v_J^2(1-\mathsf P_m)/Q$, where
$\mathsf P_m=P_m(1-Q/16)$ and $Q=-\Delta_J$.  The numerator vanishes
at $Q=0$, so the quotient has
polynomial continuation there.  Its spectral values are nonnegative
because $0\leq\mathsf P_m\leq1$, and its degree is $O(m)$, giving range
$O_J(m)$.  Continue it by the literal shells
$\Gamma_u=\Gamma_u^{\rm en}=A_0^{-1}(\mathsf P_u-\mathsf P_{Lu})$,
$u=mL^j$.  After the shells with upper side at most a chosen
$m_*=mL^h$, $h\ge0$, the remaining mean-free covariance is
$\Gamma_*^{\rm en}=A_0^{-1}\mathsf P_{m_*}$, together with the
uniform common mean.  Its positive-order terminal bounds are asserted
when $n\ge2m_*$; the activity construction separately uses the fitting
grids and their specified stopping clearance.
Their sum on nonconstant modes is exactly $A_0^{-1}$.
Polynomial constant-mode continuations are included in the independent
high integrations and are absorbed by the uniform mean.
Increasing $s$ makes $m$ exceed the fixed-$J$ starting scale in
Lemma~\ref{wh:kernels}.  For an allowed split continuation its exact multiplier is
$\lambda_j^{\rm en}$, and its balanced factor is
$\ell_j=\prod_{h\ge j}(\lambda_h^{\rm en}/\lambda)$; all projection
operators in the balance sums are those of these entry shells.
The entry-family diagonal estimate proves convergence of this product.
Average each pattern over
$\eta$ and $\zeta_{<m}$ in
$y=B_s\alpha(\psi+\zeta_{<m})+\eta$.
Choose the fixed connectivity radius to contain the covariance range
and all activity collars.  The union of covers then has connected
components whose field factors, choices, and Gaussian integrations
factor.  Grouping these components gives exactly the subset algebra
of Proposition~\ref{rg:partition-map}.  In particular it retains all exclusion
rules of the original marked expansion.  Winding links retain their
large supports.  Every operation on an unwrapped padded support is
therefore identical on plane and torus, which proves the claimed
locality once the following norm estimates have been established.

\subsection{The observation-energy regulator}

Choose, with strict margins,
\begin{equation}\label{en:reserve}
 \max(0,1-a_h/k_h)<h_1<h_*<1/p_2,
              \qquad 1<p_1<p_2.
\end{equation}
These choices are possible because $a_h/k_h>0$.  Reduce the allowed
parameter neighborhood so that its bare large-field coefficient remains
strictly below $h_1$.  For $u=mL^j$ and the enlarged site set
$D=D_u(X)$ let $E_D$ be the open observation energy, with all internal
$J$-bonds and whole $s$-cells.  Define
\begin{equation}\label{en:variational}
 V_j(X,\psi)=\frac12\sup_{\zeta,e}
 \left\{h_*E_D\bigl(B_s\alpha(\psi+\zeta)+e\bigr)
       -p_1^{-1}(\|\zeta\|_{<u}^2+\|e\|^2)\right\}.
\end{equation}
The norm $\|\zeta\|_{<u}$ is the Cameron norm of all Gaussian pieces
strictly below $u$, and $e$ has the unit-noise Cameron norm.  A singular
covariance means its Gaussian Hilbert space; thus no inverse is taken
on a null space.  Only restrictions to $D$ matter.

\begin{lemma}[Regulator composition]\label{en:regulator}
The supremum in Equation~\eqref{en:variational} is finite.  Uniformly
in scale and volume,
\begin{equation}\label{en:regulator-energy}
 0\leq V_j(X,\psi)\leq CE_D(B_s\alpha\psi)
                    \leq C\ip{\psi}{A_{0,D}\psi}.
\end{equation}
It is monotone on support enlargement and adds on compatible separated
supports.  Shell averaging satisfies the regulator estimates of
Lemma~\ref{rg:regulator} with $W^\kappa$ replaced by $W^\kappa e^{V_j}$,
and with a fixed $C^{|X|}$ in place of the former fixed per-block
constant.  For a small input polymer the half-exponent reserve in $W$
is preserved.  The same assertion holds for terminal integration.
Consequently the localization, charge estimates, and exact partition
map of Lemmas~\ref{rg:localization}, \ref{rg:charged} and
Propositions~\ref{rg:partition-map}, \ref{wh:map} retain their stated
contractions and analytic nonlinear bounds after this modification.
\end{lemma}

\begin{proof}
Write $E_D(y)=\ip{y}{My}$, and let $L$ map the Gaussian Hilbert
coordinates $\xi$ of $(\zeta,e)$ to $B_s\alpha\zeta+e$ on $D$.
The fundamental bound is
\begin{equation}\label{en:energy-contraction}
                       \|M^{1/2}L\|\leq1.
\end{equation}
For the uncut full observation covariance it is the defining inverse
covariance identity.  Restriction to an open set only lowers the
minimized energy: use the actual restricted field as trial function.
Equivalently, the field portion follows from
$A_0^{1/2}C_{<u}A_0^{1/2}\leq I$, and the observation minimization
also permits its actual noise.  The same proof includes any future
Gaussian pieces.  Constants are killed by the energy, including on
the torus.

Put $H=h_*M$ and $Q=L^*HL$.  Equation~\eqref{en:energy-contraction}
and $p_1h_*<1$ imply
\begin{equation}\label{en:optimized-form}
 2V_j=\ip{f}{Hf}
 +\ip{L^*Hf}{(p_1^{-1}I-Q)^{-1}L^*Hf},
                         \qquad f=B_s\alpha\psi.
\end{equation}
This proves finiteness, positivity, and the first bound in
Equation~\eqref{en:regulator-energy}.  The second follows by trying
$\alpha\psi$ with zero noise in the minimization defining $E_D$.
Increasing $D$ increases the minimized energy; taking the same full
Gaussian Hilbert coordinates proves monotonicity of the supremum.
For separated compatible sets the accumulated covariance has finite
range inside their fixed collars; their coordinates and energies
split, proving additivity.

For composition let $z$ be standard coordinates for the next shell
and $Tz=B_s\alpha\zeta_{\Gamma_u}$.  Completing its Gaussian square
gives the identity
\begin{align}\label{en:sup-composition}
 \left(\E_z e^{p_1V_j(f+Tz)}\right)^{1/p_1}
 &=\det(I-p_1T^*H_jT)^{-1/(2p_1)}
   \exp\left\{\frac12\sup_z
           [2V_j(f+Tz)-p_1^{-1}\|z\|^2]\right\},
\end{align}
where $H_j$ is the quadratic matrix in Equation~\eqref{en:optimized-form}.
The supremum on the right is exactly the supremum defining the next
$V$, since it combines the old coordinates and the shell coordinates
with the same Cameron penalty.  Enlarging the support only increases
it.  Applying Equation~\eqref{en:energy-contraction} to both coordinate
sets shows that the determinant has a strict stability margin depending
only on Equation~\eqref{en:reserve}.

For a positive matrix $F$ with $\|F\|\leq1-\delta$,
$-\log\det(I-F)\leq\delta^{-1}\tr F$.  The trace needed here is at
most a constant times
\[
 \E E_D(B_s\alpha\zeta_{\Gamma_u})
 \leq\E\ip{\zeta_{\Gamma_u}}{A_{0,D}\zeta_{\Gamma_u}}
 \leq C|X|.
\]
The last inequality is precisely the positive-order difference bound
for the shell kernels; the factors $u^{-2}$ cancel the number of sites
in each block.  Its constant is independent of the microscopic cell
size and of sufficiently large $L$.  The terminal covariance satisfies
the same bound.  No determinant over all microscopic height variables
is introduced by this argument.

Apply H\"older with exponents $p_1$ and $p_1/(p_1-1)$ to the two
regulators.  Choose $\kappa$ small enough for the latter power and for
the analytic reserve required by Lemma~\ref{rg:regulator}.  Its proof
then gives $W_{Lu}^{\kappa/2}$ for a small polymer after the shell
$u\to Lu$, while
Equation~\eqref{en:sup-composition} transports the factor $e^V$ without
increasing its coefficient.  Since $V_j(t\psi)=t^2V_j(\psi)$ for
$0\leq t\leq1$ and constants do not affect it, all Taylor interpolation
arguments in Lemma~\ref{rg:localization} keep this factor bounded by
its output value.  The remaining polynomial losses are absorbed by
the half-exponent reserve in $W$.  Imaginary charge shifts use only
bounded nonconstant directions, whose internal energy is $O(|X|)$;
their costs follow from the same reserve and completion.  Thus the
linear powers of $L^{-1}$ are unchanged.  The extra fixed $C^{|X|}$
is paid by the same large-component weight gain, choosing $L$ and
$A$ as in the toolkit.  Product, decoration, and regrouping estimates
use monotonicity and additivity just proved, so their quadratic
nonlinear bounds also hold.
\end{proof}

\subsection{Centered accounting and the real pattern estimate}

Fix a pattern with $j_1$ marks, $j_2$ additional selected tiles, and
selected links of lengths $d_\ell$.  All fixed factors below can depend
on the already chosen norms.  In a marked patch, prescribe its guard
heights $b_i^p$ and denote the two interior kernels by $K_o,K_c$.
Define
\begin{equation}\label{en:delta}
 \delta_i(y,b_i^p)=
 \frac{|K_o(y,b_i^p)-B_i(y)K_c(y,b_i^p)|}
 {K_o(y,b_i^p)+B_i(y)K_c(y,b_i^p)}\leq1.
\end{equation}
Factors depending only on the guard or its exterior cancel in this
ratio.  Given all guards the different marked interiors have no direct
coupling.  Therefore taking absolute values of their product of
differences gives the sum of the $2^{j_1}$ positive topologies, with
the factors $B_i$ on cut choices, times the expectation of
$\prod_i\delta_i$ in each topology.  This is an exact conditional
integration followed by an inequality; it requires no independence
of the random guards.

\begin{lemma}[Local normalization accounting]\label{en:normalization}
Let $G$ be any one of these positive topologies in a padded halo
containing all its affected tiles and interacting components.
After division by Equation~\eqref{en:base}, its contribution, apart
from the additional tile and link factors, is at most
\begin{equation}\label{en:normalization-bound}
 e^{Cj_1+Ce^{-cR}\mathcal E(y)}
 e^{E_G(y)/2}\frac{Z_G(y;k)}{Z_G(0;k)}
                \E_{G,y;k}\prod_i\delta_i(y,b_i^p).
\end{equation}
The energy $\mathcal E$ is a sum of squared nearest-cell differences
inside a fixed padded halo of the marks.  Both its size in tiles and
the number of charged centered normalizers are $O(j_1)$, uniformly in
the containing torus.
\end{lemma}

\begin{proof}
Write $\pi$ for the pattern together with a choice of its positive
topology.  After all common factors have canceled, let $\mathcal A$
be the affected halo tiles and let $\mathcal I_\pi,\mathcal I_0$ be
the remaining proxy events on this path and the no-mark path.
The set $\mathcal I_\pi$ includes every surviving unmarked replacement
and the factor $B_i$ on every marked cut choice.  Write
$\Delta E_i=E_{o_i}-E_{c_i}$, so that
$B_i=C_i e^{-\Delta E_i/2}$.  The notation $Z_G,E_G$ includes the
product and sum over all components of $G$, respectively.  Define
\[
 \mathcal Q_G^{\rm all}(y)
 =e^{E_G(y)/2}\frac{Z_G(y;k)}{Z_G(0;k)}
                  \E_{G,y;k}\prod_{i\ {\rm marked}}\delta_i.
\]
Apart from the separately selected tile factors and links, this
positive topology's contribution to the majorant after division by
Equation~\eqref{en:base} is exactly
\[
 \frac{\prod_{i\in\mathcal I_\pi}B_i^\pi(y)}
      {\prod_{i\in\mathcal I_0}B_i^0(y)}
 \frac{Z_G(y;k)\E_{G,y;k}\prod_{i\ {\rm marked}}\delta_i}
      {\prod_{b\in\mathcal A}[Z_b(0;k_h)e^{-E_b(y)/2}]}
 =\mathsf C_{\pi,G}(k)e^{\mathsf Q_{\pi,G}(y)/2}
       \mathcal Q_G^{\rm all}(y),
\]
where
\[
 \begin{split}
 \mathsf C_{\pi,G}(k)
 &=\frac{Z_G(0;k)\prod_{i\in\mathcal I_\pi}C_i^\pi}
         {\prod_{i\in\mathcal I_0}C_i^0
          \prod_{b\in\mathcal A}Z_b(0;k_h)},\\
 \mathsf Q_{\pi,G}
 &=\sum_{b\in\mathcal A}E_b-E_G
   -\sum_{i\in\mathcal I_\pi}\Delta E_i^\pi
   +\sum_{i\in\mathcal I_0}\Delta E_i^0.
 \end{split}
\]
Thus every external proxy factor is in $\mathsf C_{\pi,G}$ and
$\mathsf Q_{\pi,G}$; the defining $B_i$ inside each $\delta_i$
remains there.  The two telescopes below bound these two prefactors.

At $y=0$, a true restoration ratio for a positive bond form supported
on a core is sandwiched between its open-patch ratio and its ratio with
the outer guard pinned.  To see the direction, interpolate the bond
precision: the derivative of the logarithm of the centered partition
ratio is the negative centered expectation of the added quadratic
form.  Added exterior precision or centered pins decrease this
expectation by Lemma~\ref{st:finite-gaussian}.  The guard-pinned ratio is
independent of the exterior.  The logarithmic gap between these two
ratios is bounded by $D(H,p)$, where $H$ is the union of core tiles
and $p$ is the guard: increasing core precision as far as a full core
pin can only increase the guard pin probability further.  By
Proposition~\ref{an:pin-interaction} and the localization estimate, this
gap is bounded and at $k_h$ differs from its exponentially small
Gaussian counterpart by an arbitrarily small amount for large $s$.
Fixed-patch continuity permits Equation~\eqref{en:replacement} to
keep its constants at $k_h$ when $k$ is nearby.

To count only local errors, take the union of tiles within distance
$S_0$ of all marks, with $S_0$ larger than all skip and dependence
radii, and telescope on its enlargement through distance $2S_0$.
Perform separately the ordinary sequence ending at $G$ and the
sequence ending at the topology with no marks.  Count every event
touching this enlarged union, including events clipped at its artificial
boundary; there are $O(j_1)$ of them.  At a clipped event use the
open-patch proxy clipped in exactly the same way on both paths.
The two paths are identical near the artificial boundary, so these
proxy factors cancel in their quotient.  Their final tile integrals
also cancel off the affected halo.  The remaining true/proxy ratios
are bounded by the centered comparison just proved, one per event.
The only $k$ versus $k_h$ isolated-tile constants are those in the
affected halo; they cost a bounded factor per tile.  This gives
$\mathsf C_{\pi,G}(k)\le e^{Cj_1}$ and excludes any extensive
normalization loss.

Apply the identical double telescoping to the Gaussian energies.
Proposition~\ref{an:gaussian-localization} bounds each local-versus-true
error matrix by an exponentially small prefactor, with decay in both
cell separation and distance from that event's core.  Each matrix has
zero row sums on its unpinned components.  Express it in pair
differences and assign shortest face-cell paths inside the padded
union.  The absence of opposite-only junctions makes these paths
available.  A path of length $d$ using a fixed edge has endpoints
within its $d$-ball; polynomial ball growth and exponential decay
bound the total charge on that edge.  Summing event locations retains
$|\mathsf Q_{\pi,G}(y)|/2\le Ce^{-cR}\mathcal E(y)$, with a weaker $c$.
On a small torus the same
argument is performed on the torus itself; the number of halo tiles
is still $O(j_1)$, and image paths are charged with their full lengths.
This proves Equation~\eqref{en:normalization-bound}.
\end{proof}

Cut $G$ into boxes of side $P=R^5$ in cell units, with cuts on a
translate of the $R$ grid.  Delete every bond crossing a cut, even
when it wraps between two copies of the same box.  Retain a mark
only if its guard and the larger separating patches used for testing
that guard lie deep inside one box.  Averaging over the finitely many
translates loses at most $C(R/P)j_1$ marks.  The Gaussian energy
increase in undoing the cuts, averaged over the same translates, is
at most $C(R/P)\mathcal E(y)$.  Indeed the cut-error matrices decay
exponentially away from the cut seams; assign their pair differences
to paths and then average the probability that the seam is near a
path edge.  This probability is $O(R/P)$.  Applying Markov's
Inequality to the two nonnegative averaged quantities chooses a
translate retaining, for large $R$, at least $j_1/2$ marks while
paying the stated energy with a changed constant.  All these choices
remain inside the original padded covers.  More precisely, a ``box''
in this argument is the clipping $G\cap Q$ by a translated $P$-square
$Q$, component by component, not the full square $Q$.  No cells outside
the already chosen padded halo are added.  The underlying cells of a
clip retain their face adjacency before the topology's bond deletions;
$\mathcal E_{\rm box}$ denotes the sum of squared differences over
these adjacent cells.  Selected isolated tiles use the same convention
on their existing halos.  These are permitted rectangular clippings
of the mixed geometry.  Each contains at most $CP^2$ observation cells,
and all their energies are supported in the original $D_m(X)$.
A halo has diameter $O(R)$ and therefore meets only a bounded number
of $P$-squares.  Consequently the number of occupied clipped boxes is
at most $C(j_1+j_2)$.  A clipping with no mark halo or selected tile
contributes the exact factor one and is not counted.  Selected links
are estimated on their own paths and do not create extra boxes.

There is also an exact comparison of the non-Gaussian terms across
these box cuts.  Write $J_G(y,b;k)$ for the joint unnormalized density
of the continuous observations and all retained prescribed guards.
The shifted-versus-centered comparison of Lemma~\ref{st:finite-gaussian} gives
\begin{equation}\label{en:shifted-box}
       \frac{J_G(y,b;k)}{J_G(0,0;k)}
       \leq\frac{J_{G^{\rm cut}}(y,b;k)}{J_{G^{\rm cut}}(0,0;k)}.
\end{equation}
It applies to the affine lattice of heights and Gaussian noise with
these prescribed observations: restoring the cut bonds adds a
positive quadratic form independent of the prescribed values.
Continuous observation coordinates cause no change: fixing them
embeds the height lattice as $(h,y-B_sh)$ in height-plus-noise space.
Prescribed guards restrict this to an affine sublattice.  The parity
proof in Lemma~\ref{st:finite-gaussian} uses only sums and differences
of lattice points, and therefore also proves the shifted ratio
comparison for this lower-rank lattice, including a shift normal to
its span.  Thus no density approximation at a rare guard event is used.

Let $S$ be the event that all retained guards vanish, let
$\delta(b)=\prod_{i\ {\rm retained}}\delta_i(y,b_i^p)$, and let
$P_{\tau,0;k}$ be
the unpinned centered high law in topology $\tau$.  Dropping the
unretained factors using $\delta_i\le1$ replaces
$\mathcal Q_G^{\rm all}$ by the no-smaller factor $\mathcal Q_G$ below.
The factor
to which the comparison is applied is exactly
\[
 \begin{split}
 \mathcal Q_G(y)
 &=e^{E_G(y)/2}
       \frac{\sum_bJ_G(y,b;k)\delta(b)}{Z_G(0;k)}\\
 &=e^{E_G(y)/2}\frac{Z_G(y;k)}{Z_G(0;k)}
                     \E_{G,y;k}\delta(b).
 \end{split}
\]
Since $J_\tau(0,0;k)/Z_\tau(0;k)=P_{\tau,0;k}(S)$,
summing Equation~\eqref{en:shifted-box} gives
\[
 \mathcal Q_G(y)\le
 e^{(E_G-E_{G^{\rm cut}})(y)/2}
 \frac{P_{G,0;k}(S)}{P_{G^{\rm cut},0;k}(S)}
 \mathcal Q_{G^{\rm cut}}(y).
\]
Cutting all crossing bonds factors the last quantity over the
translated $P$-squares that meet $G$.  For each square let
$\mathfrak b$ be its entire clipped union of pieces of $G$, grouping
all current-topology components in that square in one factor.
Retention puts every retained mark's complete guard and testing
data in one such $\mathfrak b$.  Thus
\[
 \mathcal Q_{G^{\rm cut}}(y)=
 \prod_{\mathfrak b}
 \left[
 e^{E_{\mathfrak b}(y_{\mathfrak b})/2}
 \frac{Z_{\mathfrak b}(y_{\mathfrak b};k)}{Z_{\mathfrak b}(0;k)}
 \E_{\mathfrak b,y_{\mathfrak b};k}
       \prod_{i\in\mathfrak b}\delta_i(y,b_i^p)
 \right].
\]
Every denominator here is the centered, unpinned
$Z_{\mathfrak b}(0;k)$.  The expectation is under the unpinned
tilted high law of that box, with its guard data jointly random.

For the remaining centered factor, enumerate the retained guards
$p_i$, write $q_i$ for the outer guard of the enlarged testing patch,
and put $S_{i-1}=\bigcap_{r<i}\{h_{p_r}=0\}$.  The skip and retention
distances place both the cut seams and all previous $p_r$ outside
this enlarged patch.  Inside it the two current topologies therefore
have the same open graph $U_i$.  Its centered law before either
displayed pin has no hard pins.  Deleting exterior precision gives
the lower bound below; adding the centered outer pin gives the upper
bound, because that pin then separates the exterior.  Thus for
$\tau=G,G^{\rm cut}$,
\[
 P_{U_i,0;k}(h_{p_i}=0)
 \le P_{\tau,0;k}(h_{p_i}=0\mid S_{i-1})
 \le P_{U_i,0;k}(h_{p_i}=0\mid h_{q_i}=0).
\]
The log ratio of the two outer terms is
$D_{U_i,k}(p_i,q_i)$, the interaction in
Equation~\eqref{an:pin-interaction-definition} for this unpinned
local law at $k$.  Consequently
\[
 \frac{P_{G,0;k}(S)}{P_{G^{\rm cut},0;k}(S)}
 \le \prod_i
 \frac{P_{U_i,0;k}(h_{p_i}=0\mid h_{q_i}=0)}
      {P_{U_i,0;k}(h_{p_i}=0)}
 =e^{\sum_iD_{U_i,k}(p_i,q_i)}
 \le e^{Cj_1}.
\]
For the last bound Proposition~\ref{an:pin-interaction} is applied
at $k_h$ to these unpinned local graphs, and
Proposition~\ref{an:gaussian-localization} bounds their limits.
After $s$ is fixed, finite-graph continuity gives the same bound
for nearby $k$.  Combining this with the chosen seam-energy bound
therefore bounds Equation~\eqref{en:normalization-bound} by the
displayed product of centered-normalized box expressions and
\begin{equation}\label{en:seam-cost}
                     e^{Cj_1+C(R/P)\mathcal E(y)}.
\end{equation}
Conditional kernels at a retained mark are unchanged because its
guard still separates it from the box exterior.

\begin{lemma}[Good-box gain]\label{en:good-box}
Fix any finite threshold $U$.  At $(k_h,a_h)$, after $s$ is sufficiently
large, and then throughout a sufficiently small parameter neighborhood,
each box expression just obtained satisfies
\begin{equation}\label{en:good-box-bound}
 C\exp\{-cRj_{\rm deep}+R^{-12}\mathcal E_{\rm box}(y)\}
\end{equation}
whenever $\mathcal E_{\rm box}(y)\leq U^2$.  Here $j_{\rm deep}$ is
its number of retained marks.  The positive constants $c,C$ can be
chosen before $U$ and $s$.  An additional isolated selected tile has
the same small factor $e^{-cR}$ in its bounded-gradient region.
\end{lemma}

\begin{proof}
Translate each connected underlying component by a common height
period so that one observed cell value lies in $[-\pi,\pi]$.
The ratios in question are exactly invariant under this simultaneous
translation of heights and observations.  Paths in a box bound every
remaining $|y_x|$ by a polynomial in $P$ times
$1+\sqrt{\mathcal E_{\rm box}(y)}$: a spanning tree in each
connected clipped component has at most $CP^2$ cells and gives such
paths without any convexity assumption.  Thus the stipulated gradient
bound leaves a compact set of observation vectors.  There are finitely
many box geometries, positive topologies, mark choices, and translations
at fixed $R,P$.

In the Gaussian limit the centered-normalized partition factor
$e^{E_G(y)/2}Z_G(y;k_h)/Z_G(0;k_h)$ tends to one.  For a deep mark,
decompose the log ratio in Equation~\eqref{en:delta} into its centered
$y=b=0$ value, the change in its conditional centered normalization
as the guard changes from zero to $b$, and the conditional log Laplace
change due to $y$, minus its comparison-energy change.  The first
term is $O(e^{-cR})$ by Propositions~\ref{an:pin-interaction} and
\ref{an:gaussian-localization}.  The second is between
$\log r_H(b)$ and zero by
Proposition~\ref{an:conditional-pins}, with $H$ containing the restored
bonds.

For the quantitative limiting bounds let $u_c$ be the cut Gaussian
harmonic prediction of the guard data, and let $\chi_H$ equal one
near $H$ with support disjoint from the guard.  The localization
lemma gives, for every $r\geq2$,
\begin{equation}\label{en:prediction}
 \big\|\|\chi_Hu_c\|_T\big\|_{L^r}
 \leq C\sqrt r\,\operatorname{poly}(P)
 e^{-cR}\bigl(1+\sqrt{\mathcal E_{\rm box}(y)}\bigr).
\end{equation}
Here $T$ is the continuum gradient-plus-unit-cell-observation form.
The bound holds when the data have either topology's centered law:
restoring bonds decreases their covariance.  In the energy Hilbert
space with norm $\|\cdot\|_T$, write the centered random prediction as
$\mathsf A g$, where $g$ denotes standard Gaussian energy coordinates.
The localized prediction map satisfies
$\|\mathsf A\|_{\mathrm{HS}}\leq
C\operatorname{poly}(P)e^{-cR}$ by
Proposition~\ref{an:gaussian-localization}.  This normalization includes
the coefficient $1/a_h$ in $T$; $g$ has identity covariance in these
energy coordinates.  Diagonalize $\mathsf A\mathsf A^*$, with
eigenvalues $\lambda_i\geq0$.  Minkowski's Inequality in $L^{r/2}$
and the standard real Gaussian moment bound give
\[
 \big\|\|\mathsf A g\|_T\big\|_{L^r}^{\,2}
 =\Big\|\sum_i\lambda_i g_i^2\Big\|_{L^{r/2}}
 \leq\sum_i\lambda_i\|g_i\|_{L^r}^{\,2}
 \leq Cr\,\tr(\mathsf A\mathsf A^*).
\]
Prove this first for finite projections; the Hilbert--Schmidt bound
and monotone convergence give the same estimate for the full
prediction.  Covariance ordering only decreases the trace for restored
data.  A tilt by $y$ adds a deterministic prediction $\mu_y$.
The localized operator bound and the period recentering already used
on the box give
\[
 \|\mu_y\|_T\leq C\operatorname{poly}(P)e^{-cR}\|y\|_2
 \leq C\operatorname{poly}(P)e^{-cR}
              \bigl(1+\sqrt{\mathcal E_{\rm box}(y)}\bigr),
\]
enlarging the fixed polynomial if necessary.  The triangle inequality
proves Equation~\eqref{en:prediction} with the displayed $\sqrt r$
dependence.  Its constant and polynomial degree are independent of
$r$, so the subsequent choice $r=O(j_{\rm deep})$ has only polynomial
moment losses.

The negative log Gaussian full-core pin ratio is at most
$\|\chi_Hu_c\|_T^2/2$: this cutoff corrects the conditional mean
to zero on $H$ without changing its guard values.  The linear-in-data
part of the conditional Laplace discrepancy is bounded by
$C\|y\|\|\chi_Hu_c\|_T$.  To verify this also when restoring bonds
glues two free faces, first prescribe a finite list of smooth guard
probes.  The cut and restored predictions minimize the same cut
energy on nested admissible spaces.  Subtracting $\chi_Hu_c$ from
the cut minimizer leaves all guard probes unchanged and makes it
admissible in the restored, trace-matching subspace.  Orthogonal
projection then bounds the difference of minimizers by
$\|\chi_Hu_c\|_T$.  The conditional quadratic parts differ from the
open-patch comparison in Equation~\eqref{en:replacement} by at most
$Ce^{-cR}\|y\|^2$, again by guarded localization.  All these arguments
use finite-energy predictions, so they never prescribe pointwise
values of a continuum rough Gaussian trace.

To transfer these estimates to the microscopic guards, first work
under the centered high laws and use finite smooth guard lists and
finite smooth core penalties.  Proposition~\ref{an:conditional-pins}
gives the joint conditional convergence for each topology in its
own centered law.  Increasing the probe lists and the core penalty
approximates the full pin ratio there: the centered identity
$\E r_H=e^{-D(H,p)}$ and Proposition~\ref{an:pin-interaction}
close the approximation error, and the Gaussian ratio is strictly
positive almost surely.

For a cut centered law $\mu$ and a restored centered law $\nu$
with identical exterior, let $A_s$ be any failure event about their
common guard data.  Equations~\eqref{an:topology-density}
and~\eqref{an:pin-ratio-away-zero} give
\[
 \nu_p(A_s)\le e^{D_U(H,p)}\mu_p(A_s),\qquad
 \mu_p(A_s)\le\mu_p(r_H\le\varepsilon)
                    +\varepsilon^{-1}\nu_p(A_s).
\]
The interactions have finite limits.  Taking $s\to\infty$ and
then $\varepsilon\downarrow0$ transfers vanishing failure
probabilities in both directions.  Thus both topologies'
conditional approximations hold jointly when their guard data are
sampled under either centered law.

Only after this centered transfer, change the sampling law separately
in each topology $\tau$.  For a fixed observation vector $y$,
\[
 \frac{\dd P_{\tau,y;k_h}}{\dd P_{\tau,0;k_h}}(h)
 =\frac{e^{(B_sh,y)}}{\E_{\tau,0;k_h}e^{(B_sh,y)}} .
\]
At fixed $R,P,U$ these are tilts by a fixed finite list of cell
observations with bounded bare covariance and coefficients in a
compact set.  The displayed densities have a common $L^q$ bound
for every fixed finite $q$, by the density moment estimate in
Proposition~\ref{an:conditional-pins}.  H\"older in each topology
therefore transfers the centered failure bounds to its own tilted
law.  There are only finitely many topologies and required moment
orders at these fixed parameters, so the choices are common to all
of them.  On exceptional data use $\delta_i\leq1$.
Consequently bounded moments of $\delta_i$ have the Gaussian upper
bounds derived from Equation~\eqref{en:prediction}; no uniform claim
for arbitrary prescribed microscopic guards is needed.

Use $|(x-1)/(x+1)|\leq|\log x|$ for $x>0$.  H\"older over the
$j_{\rm deep}\leq\operatorname{poly}(P)$ retained marks gives their
product a factor $e^{-cRj_{\rm deep}}$, with polynomial losses in
$P$, the moment order, and $1+\mathcal E_{\rm box}$.  For clarity,
if a per-mark bound has the form
$e^{-c_0R}P^Cj^C(1+\mathcal E)^C$, maximize
$(1+\mathcal E)^{Cj}e^{-R^{-12}\mathcal E}$ to obtain a factor
$[C(1+jR^{12})^C]^j$.  Its logarithm is $O(j\log R)$ because
$j\leq\operatorname{poly}(P)$, and is absorbed by $c_0Rj/2$.
This proves Equation~\eqref{en:good-box-bound} in the limiting
Gaussian calculation with slack.

The same argument works along every convergent sequence in the compact
set of translated observation vectors, using uniform exponential
integrability of the finite probes.  A subsequence contradiction
therefore makes the bound uniform there for sufficiently large $s$.
All finitely many box calculations are then continuous in $k,a'$;
their strict slack gives a common neighborhood.  This last continuity
is on fixed finite microscopic graphs.  It invokes no scaling theorem
away from $k_h$.  Equation~\eqref{en:tile-error} is handled by the
free mixed-domain convergence in exactly its own finite tile: its
normalized ratio tends to one uniformly on the same compact sets,
so its error can be made smaller than $e^{-cR}$.
\end{proof}

In a box outside this gradient threshold discard its $\delta_i$'s.
Completing the square and centered Gaussian domination give
\begin{equation}\label{en:bad-box}
 e^{E_{G_b}(y)/2}\frac{Z_{G_b}(y;k)}{Z_{G_b}(0;k)}
 \leq\exp\left\{\frac12\max(0,1-a'/k)E_{G_b}(y)\right\}.
\end{equation}
Indeed the bare comparison observation covariance with coefficient
$k$ is bounded by $\max(1,k/a')$ times the reference covariance;
inverting on the complement of the constant mode gives the stated
energy coefficient.  Constants cancel on that mode.  Fixed centered
normalization ratios at nearby parameters cost a bounded factor per
active tile.  For a bad selected isolated tile, the bound for
$|F_b|$ follows by adding one to Equation~\eqref{en:bad-box}.
Group these tiles also into boxes of side at most $P$.

On every bad box insert the upper bound
\begin{equation}\label{en:threshold}
       1\leq\exp\{\varepsilon_1(
                \mathcal E_b(y)-U_{\mathrm{big}}^2)\},
\end{equation}
where $\varepsilon_1>0$ is fixed sufficiently small.  Differences are
always those on the underlying uncut cells, even if the current
topology has fewer bonds.  Choices of good and bad boxes cost only
an exponential per active tile.  The real estimate now consists of
the mark and tile gains in the good boxes, the tiny seam cost
$C(R/P)\mathcal E$, the link costs in Equation~\eqref{en:links},
and the broad quadratic costs in bad boxes with the negative constants
from Equation~\eqref{en:threshold}.

\subsection{Averaging without a microscopic volume loss}

Let $X$ be the padded $m$-block cover of the pattern and $D=D_m(X)$.
All its positive quadratic costs are bounded by $h_1E_D/2$ once
$R$ is large and $\varepsilon_1$ is small.  Here are the two required
form comparisons.  First, the sum of cut-box minimized energies is
at most $E_D$: restrict any trial field in the latter minimization to
the boxes and discard crossing bonds.  Second, on two neighboring
observation cells, cell Poincar\'e and the observation penalty give
\[
 |y_x-y_{x'}|^2\leq C\left[
  \sum_{z\in\{x,x'\}}|(B_sv)_z-y_z|^2
             +(v,A_{0,D}v)/a'\text{ on these cells}\right].
\]
Sum this inequality over the good-box edges with coefficient
$R^{-12}$, and over the localized seam and path edges with their
coefficients.  Their multiplicities are bounded: the boxes are disjoint
clippings and the selected-tile halos have fixed overlap, while the
seam and path bounds were proved above.  Minimizing over $v$ gives
$\sum_{\rm good}R^{-12}\mathcal E_{\rm box}\le CR^{-12}E_D$
and bounds the seam and path total by another small multiple of $E_D$.
Equation~\eqref{en:reserve} supplies strict room beyond the broad
coefficient in Equation~\eqref{en:bad-box}.

Gaussian completion bounds the resulting background quadratic by
$V_0$ in Equation~\eqref{en:variational}.  The determinant actually
integrated is estimated with its localized costs rather than with
the larger form $h_1E_D$.  Its covariance contraction has a strict
margin from one.  Its trace is bounded by
\begin{equation}\label{en:trace-entry}
 C\operatorname{poly}(P)N_{\mathrm{bad}}
 +CR^{-12}P^2N_{\mathrm{good}}
 +C(R/P)R^2(j_1+j_2)
 +C\sum_\ell e^{-c(R+d_\ell)}\operatorname{poly}(d_\ell).
\end{equation}
Here $N_{\mathrm{bad}}$ and $N_{\mathrm{good}}$ count occupied
clipped boxes, not all boxes in the containing torus.  For the first
term the rank is the number of observed cells in the bad boxes, at
most $CP^2$ per box, and the covariance contraction bound controls
every eigenvalue.  For the good-box term, every adjacent-cell
difference of $B_s\alpha\zeta_{<m}+\eta$ has variance bounded by a
constant independent of $s,m$: the height part is bounded by the
nearest-cell flow comparison and the two independent noise variables
have variance two.  There are at most $CP^2$ such differences per
clipping.  Multiplication by the allowance $R^{-12}$ therefore gives
$CR^{-12}P^2$ per occupied good box.  This accounts for the whole
observation rank; it uses no cancellation of its noise coordinates.
For the seam and link terms sum their localized coefficients against
the same adjacent-cell variance bound.  There are $O(R^2)$ cells per
halo item.  The strict covariance margin converts this trace estimate
into the corresponding logarithmic determinant bound.

Thus Equation~\eqref{en:trace-entry} is independent of $s$: the
determinant never counts the microscopic height dimension.  Since
$P=R^5$ and $N_{\mathrm{good}}\le C(j_1+j_2)$, both the good-box
allowance and the seam term are $O(R^{-2})(j_1+j_2)$, absorbed by
the real per-item gain.  This is the reason for choosing the allowance
$R^{-12}$ in Lemma~\ref{en:good-box}.  Its polynomial maximization
cost is still only $O(j_{\rm deep}\log R)$, so that lemma's constants
remain independent of the later threshold and of $s$.

A bad box contains at most a polynomial number of marks and tile
items.  Its loss of their gains is at most a polynomial in $P,R$
in the logarithm.  Choose the finite $U_{\mathrm{big}}$ so that
$\varepsilon_1U_{\mathrm{big}}^2$ exceeds this loss, the first term of
Equation~\eqref{en:trace-entry}, and all per-box combinatorial costs,
leaving a further factor $e^{-cR}$ per item.  Only after this choice
apply Lemma~\ref{en:good-box} to choose $s$.  The finite choice of
optimized cut translate does not create a random determinant: replace
the pointwise optimum by a sum over its at most polynomially many
translates.  For a nonempty configuration this polynomial is absorbed
by the per-item exponential gain.  The all-empty configuration has
no cut choice or cost.

We obtain for the Gaussian-averaged pattern $F_\pi$ the real estimate
\begin{equation}\label{en:averaged-pattern}
 |F_\pi(\psi)|\leq
 C^{j_1+j_2}\exp\{-cR(j_1+j_2)\}
 \prod_\ell Ce^{-c(R+d_\ell)}
 W_m^{\kappa_0}(X,\psi)e^{V_0(X,\psi)},
\end{equation}
with $\kappa_0<\kappa$ and fixed reserve in both quadratic estimates.
The optional $W^{\kappa_0}$ records the small reserved path and
direction costs.  Completion also permits replacing the background
by $\psi+xf$, for a real direction of scaled point norm at most one,
at extra cost
\begin{equation}\label{en:real-translate}
                              e^{C x^2|X|}.
\end{equation}
To see the bound with the regulator still evaluated at $\psi$, use
the strict gap $h_1<h_*$ and Cauchy--Schwarz in the optimized quadratic
form.  The excess on the direction is bounded by
$C\ip{f}{A_{0,D}f}\leq C|X|$, since its first differences are
$O(m^{-1})$ and there are $O(m^2|X|)$ sites.  The bounded-range
corner and collar bonds change only this fixed constant.  This
reserve is also why a translate does not consume the final regulator.

Equation~\eqref{en:averaged-pattern} proves absolute integrated
convergence for the link expansion.  For a finite torus first sum the
finite positive topologies for each microscopic height configuration,
then perform the Gaussian integrations just bounded, and finally sum
links.  Tonelli for absolute values and the resulting summability
justify all exchanges in the signed identity.

\subsection{Complex fields, analytic directions, and input derivatives}

A separate bound on imaginary directions is necessary because taking
absolute values inside each conditionally oscillating guard kernel
would lose the cancellation in its Gaussian energy.  We give the
finite theta identity that avoids that loss.  For a positive definite
precision $T$ on the unpinned height coordinates set
\[
 w(t)=\sum_{h\in\Lambda}e^{-\ip h{Th}/2+\ip h t},\qquad
 w_c(t)=\sum_{p\in\Lambda+c/2}e^{-\ip p{Tp}+\ip p t},
                 \quad c\in\Lambda/2\Lambda.
\]
Sum and difference of the two lattice variables give, for real $x,z$,
\begin{equation}\label{en:theta-parity}
 w(x)w(iz)=\sum_c|w_c(x+iz)|^2,
 \qquad w(0)w(x+iz)=\sum_cw_c(x+iz)^2.
\end{equation}
The cosets are invariant under sign reversal, which explains the
second equality.  The triangle inequality proves
$|w(x+iz)|w(0)\leq w(x)w(iz)$.  Since
$Z(y)=e^{-\|y\|^2/2}w(B_s^*y)$, including its Gaussian prefactors
gives the useful form
\begin{equation}\label{en:complex-comparison}
             \frac{|Z(y+iv;k)|}{Z(y;k)}
                         \leq\frac{Z(iv;k)}{Z(0;k)}.
\end{equation}
Every free component's constant is controlled by the observation mass.
Hard-pinned subspaces use their own lattice, so the identity applies
there too; a vanishing auxiliary mass also deals with any redundant
semidefinite presentation.

The centered imaginary ratio equals
$e^{\|v\|^2/2}\E e^{i\ip{B_sh}v}$, whose characteristic factor is
positive.  Centered zero pinning increases this characteristic
function.  In detail, condition on the pinned coordinates; the
remaining law is a real tilt of the centered pinned Gaussian lattice.
Equation~\eqref{en:complex-comparison} bounds its characteristic
modulus by that of the centered pinned law.  Average the conditional
characteristic functions and use positivity of the centered full
characteristic function.  This proves the claimed direction without
taking absolute values of individual lattice Fourier summands.

For an imaginary observation direction $v$ of any prescribed fixed
supremum, expand each microscopic difference into its positive
topologies, discarding difference rarity for the moment but retaining
the link gains.  Apply Equation~\eqref{en:complex-comparison} to each
whole integrated topology.  Pin whole strips of width $R/4$ along
translates of the $P$-box cuts.  These thick walls separate the region
into bounded boxes.  Bonds can then be cut through the middles of the
zero strips, since every deleted endpoint is already zero.  All strips
are intersected with the same active halo; neither walls nor auxiliary
boxes enlarge its support.  Thus each pre-pin component is a permitted rectangular clipping, exactly as
required for Proposition~\ref{an:mixed-scaling}.

On these fixed pinned boxes that lemma gives the Gaussian imaginary
ratio, uniformly for bounded $v$, with arbitrarily small multiplicative
error per box after $s$ is large.  It remains true nearby by finite-graph
continuity.  The additional Gaussian pin energy over the unpinned
energy is, after averaging wall translates, at most
\begin{equation}\label{en:imaginary-wall}
       C\|v\|_\infty^2(R/P)\,\#\{\text{active cells}\}.
\end{equation}
For this bound use the massive-form localization away from each wall;
near it use the uniform observation-energy bound.  The proportion
of near-wall cells is $O(R/P)$ and the decaying tails have the same
average.  There is no zero-row-sum requirement here because $v$ is
bounded.  The optimizing wall translate need not equal the translate
used for the real estimate.

The Gaussian shape factors in Equation~\eqref{en:base} cancel this
Gaussian imaginary ratio: for a quadratic $E$,
$\Re E(y+iv)=E(y)-E(v)$.  The double telescoping in
Lemma~\ref{en:normalization} bounds the remaining shape discrepancies
by exponentially small localized costs.  For the real part of each positive topology we use the following
separate consequence of the good-box proof, with no $\delta_i$ factors:
\[
 e^{E_{G_b}(y)/2}\frac{Z_{G_b}(y;k)}{Z_{G_b}(0;k)}
       \le 2\exp\{R^{-12}\mathcal E_b(y)\}
       \qquad(\mathcal E_b(y)\le U_{\mathrm{big}}^2).
\]
At $(k_h,a_h)$ its left side tends uniformly to one on the compact
set of recentered observations, by the same fixed mixed-domain
convergence.  Choose $s$ so that it is at most $3/2$ for every one
of the finitely many positive topologies.  Finite-graph continuity
then makes it at most $2$ on one parameter neighborhood.  Thus the
displayed numerical constant is independent of $R,P$ and the threshold;
only the required $s$ and neighborhood depend on these earlier choices.  This establishes the estimate without deleting
factors from an inequality that contained $\prod_i\delta_i$.
For bad boxes use Equation~\eqref{en:bad-box} and the same threshold
penalty.  Hence the real-part determinants are precisely the localized
determinants already bounded.
In particular the good-box allowance still costs at most
$CR^{-12}P^2N_{\mathrm{good}}\le CR^{-2}(j_1+j_2)$, even though
no mark gain is used on this strip boundary.  Bad-box determinants
are paid by their threshold penalties as before, and empty boxes
remain exact identities.  Thus this cost belongs to the $o(R)$
term below.  Links at complex arguments retain their tiny path costs directly.
For a fixed strip height $H$ we have therefore proved the second bound
\begin{equation}\label{en:strip-rough}
 \begin{split}
 |F_\pi(\psi+(x+iy)f)|\leq{}&
 \exp\{o(R)(j_1+j_2)+C_H(1+x^2)|X|\}
 \prod_\ell Ce^{-c(R+d_\ell)}\\
 &\hspace{5mm}\cdot W_m^\kappa(X,\psi)e^{V_0(X,\psi)},
                         \qquad |y|\leq H.
 \end{split}
\end{equation}
Here $H$ is fixed before $R$; the $o(R)$ includes
Equation~\eqref{en:imaginary-wall}, since each halo has $O(R^2)$ cells
and $R^3/P=R^{-2}$.  On each fixed finite calculation analyticity
follows from Gaussian domination with quadratic reserve.  The same
domination and the absolute link sum prove analyticity after averaging.

Choose a desired radius $r_0>4eh_0$ and then $H>4r_0$.
For fixed real $\psi,f$, apply the three-lines theorem on each of the
strips $0\leq\pm\Im z\leq H$ to
$F_\pi(\psi+zf)e^{-D z^2|X|}$, with $D>C_H$.
Equation~\eqref{en:real-translate} bounds its real boundary and
Equation~\eqref{en:strip-rough} its other boundary.  On $|z|\leq r_0$
the real-boundary harmonic weight is at least $3/4$.  Consequently
the gain $e^{-cR(j_1+j_2)}$ survives with a smaller $c$; the damping
and its removal cost only $e^{C_{r_0,H}|X|}$.  Moreover
\begin{equation}\label{en:cover-size}
       |X|\leq C(j_1+j_2)+C\sum_\ell(1+d_\ell/R),
\end{equation}
so these losses are absorbed by the item and link gains after increasing
$R$.  We have recovered Equation~\eqref{en:averaged-pattern}, with
weaker constants, on the complex directional disk of radius $r_0$.

Cauchy's Inequality bounds the diagonal $r$th derivative by
$r!r_0^{-r}$ times this bound.  Real polarization bounds the
multilinear derivative by a further factor $r^r/r!$.  Thus its
contribution to the point norm is at most
$(h_0/r_0)^r r^r/r!$ times the pattern bound, and the sum over $r$
converges because $r_0>4eh_0$.  This proves the full analytic point
norm estimate, including independent real directions in each slot.

For input derivatives, fix the already chosen $s,R$.  Every scalar
denominator in the unsummed expansion is independent of $k$, by
Equation~\eqref{en:replacement} and Equation~\eqref{en:base}.
On a sufficiently small complex $k$-disk about $k_h$, absolute values
of the height weights are bounded by weights at a nearby real
precision.  Discarding the imaginary-direction cancellation and
using broad bare quadratic completion on every active tile now costs
at most $e^{C_{s,R}(j_1+j_2)}$ per pattern.  This constant is finite:
there are finitely many microscopic sites per active halo, their
centered moments are bounded under the observation mass, and the
local normalization accounting compares only these halos.  The
quadratic stability margin and exponential link decay persist on
a smaller disk.

On its real diameter we already have the small complex-field estimate.
Apply the two-constants principle separately to the upper and lower
half-disks, after the same field-translation damping.  On a concentric
disk of relative radius $\rho$ the harmonic measure of the circular
boundary is $O(\rho)$.  Choose $\rho$ so small that its contribution
$O(\rho C_{s,R})$ is less than a fixed fraction of the remaining
$cR$ per item.  The pattern estimate is then still exponentially small
on that complex $k$-neighborhood.  Cauchy's Inequality there proves
Equation~\eqref{en:derivatives}; for example one may take constants
bounded by $r!\rho_k^{-r}$ times a fixed factor on an inner disk.
This interpolation is performed on each local pattern integral, not
on a volume-dependent total partition function.  Finally, varying
$a'$ changes finite Gaussian integrals, the quadratic forms, and the
common frequency continuously.  The same absolute estimates with
reserve give dominated convergence on each fixed support, for each
field derivative and scalar projection.  This proves the stated
continuity in $a'$ without asserting differentiability of a norm
whose period changes.

\subsection{Connected pattern sums and completion of entry}

After the preceding estimates assign a positive weight $w(I)$ to each
halo item and link, including all its analytic and geometric constants.
For a fixed $m$-block $B$, their placement sums obey
\begin{equation}\label{en:placement}
 \sup_B\sum_{I:\,\operatorname{cover}(I)\text{ meets the
       connection neighborhood of }B}
       w(I)(2A)^{|\operatorname{cover}(I)|}
                         e^{|\operatorname{cover}(I)|}
                 \leq e^{-cR}=:\eta_R.
\end{equation}
Indeed a halo has bounded block diameter and only polynomially many
cell-root and event labels in a fixed $m$-block.  A link of length $d$
uses at most $C(1+d/R)$ blocks.  Its endpoints and path descriptions
have polynomially many choices at each fixed length; allowing
translations and cover attachments adds at most an exponential in
$1+d/R$.  These factors, the weight $(2A)^{|\operatorname{cover}|}$,
and the displayed extra exponential are absorbed by
$e^{-c(R+d)}$.  All counted cells are observation cells; their count
is independent of $s$.  Repeated roots use subsets of their finite
local labels, bounded by the corresponding product of $1+w(I)$.

Here is a direct connected-sum bound.  Root a connected collection
of items at an item meeting a specified block, and choose a canonical
spanning tree of the item adjacency graph.  A child can attach to any
block of its parent's cover.  At bounded depth let the sum of all
possible descendant trees attached at a fixed block be $T$.  The
unordered-child sum is bounded by its exponential generating bound,
so for a parent cover $S$ it is at most $e^{|S|T}$.  Starting with
depth zero and using Equation~\eqref{en:placement}, induction gives
$T\leq\eta_R<1$ at every depth: the factor $e^{|S|}$ reserved there
pays all descendants.  Monotone convergence bounds the sum of all
finite trees.  Canonically assigning each actual collection to one
tree only decreases this positive overcount.  The factor $2^{|S|}$
in Equation~\eqref{en:placement} pays the choices of occupied covered
blocks and regrouping into the subset convention.  Therefore the
connected pattern sum at weight $A$ is at most $C\eta_R$ in the norm
of Equation~\eqref{en:norm}.  The same encoding works for projected
torus covers rooted at actual blocks; a winding item's full cover and
length have already been charged.

Choose $R$ large enough that $C\eta_R\leq\epsilon$, and then make
the choices of threshold, $s$, and parameter neighborhood in the order
stated in Theorem~\ref{en:entry}.  The expansion in
Equations~\eqref{en:replacement}--\eqref{en:links} is exact, the
integrated sums are absolutely convergent, and grouping separated
components is exactly the partition algebra.  This proves
Equation~\eqref{en:small} and its torus version.  Equation~\eqref{en:placement}
with the parameter derivative costs proves the differentiated bounds
after the same possible reduction of the neighborhood.

For completeness, inserting a bounded function of $Y$ in
Equation~\eqref{en:exact-density} preserves the exact observation
identity before high-mode averaging; all local expansions are identities
in its integrand.  Subsequent smooth-source tests are implemented by
their Gaussian translations, with the same normalized denominator.
Thus the entry does not assume that a height observable already equals
an observable of the auxiliary Gaussian field.  Its purpose is the
exact local interaction representation to which those translations
and the terminal diagnostics apply.  Finally, shell range and padded
supports imply plane copying through the stopping scale by
Proposition~\ref{wh:map}; large noncopying torus supports have the same
weight-gain estimates.  Lemma~\ref{en:regulator} supplies all regulator
hypotheses for continuing that map.  This completes the proof of
Theorem~\ref{en:entry}.
